\section{Original Paper and Its Weaknesses}
\label{sec:original}

\subsection{FakeShield Framework Overview}

FakeShield~\citep{fakeshield2025} is a multimodal framework for explainable image forgery detection and localization. It consists of two primary modules:

\subsubsection{Domain-Tag Enhanced Forgery Detection Module (DTE-FDM)}

DTE-FDM is built upon LLaVA-v1.5-13B~\citep{llava}, a vision-language model that combines a CLIP visual encoder with a LLaMA language model. The key innovation is the incorporation of domain-specific tags (AIGC, DeepFake, Photoshop) generated by a Domain Tag Generator (DTG)---a ResNet-50 classifier trained on three forgery categories. These tags are prepended to the input prompt, providing the model with prior knowledge about the suspected manipulation type.

The DTE-FDM takes an image and domain tag as input, and outputs:
\begin{itemize}
    \item A binary detection decision (tampered or authentic)
    \item A natural language explanation identifying specific artifacts
    \item References to similar training examples for MFLM localization
\end{itemize}

\subsubsection{Multi-modal Forgery Localization Module (MFLM)}

MFLM is based on GLaMM~\citep{glamm} and SAM~\citep{sam}, leveraging grounding capabilities to produce pixel-level segmentation masks of tampered regions. It uses the DTE-FDM's output (detection result and training data references) to guide the localization process.

\subsection{Identified Weaknesses}

Through systematic reproduction, we identified four critical weaknesses:

\subsubsection{Weakness 1: Pipeline Path Resolution Bug}

The DTE-FDM outputs absolute file paths from the training environment (e.g., \texttt{/root/hw/multi-media/FakeShield/playground/image/xxx.jpg}). When deployed in a different environment, these paths do not exist, causing MFLM to fail with file-not-found errors. In our test environment, this resulted in a 0\% end-to-end success rate before applying our fix.

\subsubsection{Weakness 2: CLIPVisionTower Loading Incompatibility}

The \texttt{load\_pretrained\_model} function in \texttt{builder.py} contains the following conditional check:

\begin{verbatim}
if 'llava' in model_name.lower():
    # Initialize vision tower
    vision_tower = model.get_vision_tower()
    vision_tower.load_model()
    image_processor = vision_tower.image_processor
\end{verbatim}

For the DTE-FDM model, \texttt{model\_name} is ``DTE-FDM'', which does not contain the substring ``llava''. Consequently, the entire vision tower initialization block is skipped, leaving \texttt{image\_processor = None}. This causes an \texttt{AttributeError} during inference when \texttt{process\_images()} attempts to access \texttt{image\_processor.image\_mean}.

This bug is particularly insidious because it only manifests when the model is loaded with a name that does not match the ``llava'' pattern, making it environment-dependent and difficult to diagnose.

\subsubsection{Weakness 3: High Memory Requirements}

The 13B-parameter DTE-FDM requires 6.94 GB of GPU VRAM in FP16 mode. This limits deployment to systems with high-end GPUs. The original paper does not explore quantization strategies for reducing memory footprint.

\subsubsection{Weakness 4: Limited Domain Tag Coverage}

The DTG classifier supports only 3 forgery categories:
\begin{enumerate}
    \item AIGC (AI-generated content editing)
    \item DeepFake (face manipulation)
    \item Photoshop (traditional splicing/copy-move)
\end{enumerate}

This taxonomy fails to capture emerging manipulation techniques:
\begin{itemize}
    \item \textbf{Stable Diffusion Inpainting}: Using SD's inpainting capabilities for seamless object insertion/removal
    \item \textbf{ControlNet-based Manipulation}: Precise structural control via ControlNet~\citep{controlnet}
    \item \textbf{Midjourney-generated Content}: High-quality AI-generated images from text prompts
\end{itemize}

As these tools become more accessible, the inability to distinguish their outputs limits FakeShield's practical utility.

\subsubsection{Weakness 5: Lack of IoU Evaluation}

The original paper evaluates MFLM qualitatively through visual inspection but lacks a systematic quantitative evaluation using Intersection over Union (IoU) metrics. Without IoU scores, it is difficult to compare localization accuracy across different methods or assess improvements objectively.
